% Generated draft: requires mathematical and source review.
\documentclass[11pt]{article}
\usepackage[margin=25mm]{geometry}
\usepackage{fontspec}
\setmainfont{Latin Modern Roman}
\usepackage{amsmath,amssymb,mathtools,unicode-math}
\setmathfont{Latin Modern Math}
\usepackage{enumitem,needspace,xurl,hyperref}
\setlength{\parindent}{0pt}
\setlength{\parskip}{5pt}
\setlength{\emergencystretch}{3em}
\newcommand{\N}{\mathbb{N}}
\newcommand{\Z}{\mathbb{Z}}
\newcommand{\Q}{\mathbb{Q}}
\newcommand{\R}{\mathbb{R}}
\newcommand{\C}{\mathbb{C}}
\newcommand{\F}{\mathbb{F}}
\newcommand{\A}{\mathbb{A}}
\newcommand{\PP}{\mathbb{P}}
\newcommand{\E}{\mathbb{E}}
\newcommand{\eps}{\varepsilon}
\newcommand{\dd}{\,\mathrm{d}}
\begin{document}
\section*{UnsolvedMath ID: 30005004}
\textbf{Source problem code:} catalogue entry
\subsection{Definitions and mathematical statement}
For an integer \(n\geq 1\), let
\[
\Pi_n=\{p\in\R[x,y]:\deg p\leq n\},\qquad N=\binom{n+2}{2}.
\]
A set \(X=\{p_1,\ldots,p_N\}\subset\R^2\) of distinct points is a \(GC_n\) set if, for every \(i\in\{1,\ldots,N\}\), there are nonconstant affine linear polynomials \(\ell_{i1},\ldots,\ell_{in}\in\R[x,y]\) such that
\[
Q_i=\prod_{k=1}^{n}\ell_{ik},\qquad
Q_i(p_j)=\delta_{ij}\quad(1\leq j\leq N),
\]
where \(\delta_{ij}=1\) for \(i=j\) and \(0\) otherwise. These are the fundamental interpolation polynomials. Here an affine linear polynomial has the form \(ax+by+c\), with \((a,b)\ne(0,0)\). For such a polynomial \(\ell\), write
\[
V(\ell)=\{(x,y)\in\R^2:\ell(x,y)=0\}.
\]
A maximal line for \(X\) means a line containing exactly \(n+1\) points of \(X\).

\textbf{Gasca--Maeztu maximal-line conjecture.} For every integer \(n\geq1\) and every \(GC_n\) set \(X\subset\R^2\), there exist an index \(i\in\{1,\ldots,N\}\) and a nonconstant affine linear polynomial \(\ell\) such that
\[
\ell\mid Q_i\quad\text{in }\R[x,y],\qquad
\lvert X\cap V(\ell)\rvert=n+1.
\]
In particular, every such set contains a maximal line. The supplied excerpt attributes this conjecture to Gasca and Maeztu in 1982.

\textit{Qualifications requiring review.} The extracted cardinality is typographically damaged; \(\binom{n+2}{2}\) is its intended binomial reading. The phrase ``linear forms'' is interpreted as affine linear polynomials defining lines in \(\R^2\). The source does not explicitly give the lower bound on \(n\); the displayed conjecture uses positive degrees. If degree zero is included, the empty-product definition gives a singleton with fundamental polynomial \(1\): the maximal-line question still has an affirmative interpretation, but the additional divisibility condition cannot hold. The accompanying inductive discussion says to adjoin \(n\) points when passing from degree \(n-1\) to degree \(n\); this conflicts with the cardinality difference \(\binom{n+2}{2}-\binom{n+1}{2}=n+1\) and is not used as a hypothesis.
\subsection{Short English statement}
The conjecture asks whether every planar interpolation set whose fundamental polynomials each factor into \(n\) affine linear factors has a line through \(n+1\) of its points. The supplied report also requires that this line divide some fundamental polynomial. Positive degrees are used here; the source's degree-zero boundary and damaged extraction require review.
\subsection{Sources}
\begin{itemize}
\item Supplied local UnsolvedMath registry record, permanent ID \(30005004\), title: Gasca--Maeztu Maximal-Line Conjecture; exact source problem code: \texttt{catalogue entry}.
\item Supplied report citation: \emph{Mini-Workshop: Interpolation, Approximation, and Algebra} (2023), Oberwolfach Reports, DOI \texttt{10.4171/owr/2022/7}, pp.\ 383--428. The extracted passage identifies Definition 1.1 and Conjecture 1.2 and includes a page marker for p.\ 385. \url{https://doi.org/10.4171/owr/2022/7}. The record separately gives proposed year 2022 and the literature-source label ``Oberwolfach Report 7/2022''; the differing year labels are retained without reconciliation.
\item The supplied excerpt attributes the conjecture to Gasca and Maeztu (1982), citing its internal reference [21]. It credits Chung and Yao with the interpolation configuration, citing [13], and mentions historical antecedents in Berzolari's 1914 work [4] and subsequent work of Radon [31]. Full bibliographic entries for these internal references are not supplied.
\item Supplied literature source: Hakopian--Vardanyan--Vardanyan, \emph{On the Gasca-Maeztu conjecture for n=6}. \url{https://arxiv.org/abs/2208.06784}.
\item Supplied literature source: Vardanyan, \emph{A new proof of the Gasca-Maeztu conjecture for n=5}. \url{https://arxiv.org/abs/2108.12799}.
\item Literature assessment dated exactly \texttt{2026-08-22}: the supplied curation labels the problem partially solved and states that the conjecture is proved for degrees \(n\leq5\), has partial progress for \(n=6\), and remains open in general. This is supplied dated curation derived from the OWR record, not an independently verified literature assessment.
\item The registry labels the statement \texttt{corrected\_verified} and reports a check against the official Oberwolfach report. This verification claim is supplied provenance; no source was independently visited or checked for this draft.
\end{itemize}
\end{document}
